\chapter{Estructura, empenaje y tren de aterrizaje}
\label{cap:mm-estructura}

\opening{Las estructuras de tubos soldados tienen alivio de tensiones: ninguna reparación por soldadura está permitida fuera de la fábrica.}

\section{Descripción de la célula}
\fuente{MM, capítulos 4 y 5}

El R22 es biplaza, de rotor principal y motor únicos, construido principalmente en metal y con tren de patines. La estructura primaria del fuselaje es de tubos de acero soldados y chapa de aluminio remachada; el cono de cola es monocasco (los revestimientos de aluminio soportan las cargas primarias). La fibra de vidrio y los termoplásticos se usan en la estructura secundaria de la cabina, los carenados de enfriamiento del motor y otros ductos y carenados. Una puerta de capó a la derecha da acceso a la caja principal y la transmisión; paneles y capós removibles dan acceso a los controles. Hay cortafuegos de acero inoxidable delante y encima del motor.

\section{Fuselaje}

La cabina es un conjunto no reemplazable en campo.

\begin{itemize}
\item \textbf{Cortafuego vertical:} su reemplazo debe hacerse en fábrica, en una plantilla. Las reparaciones pueden hacerse según la AC 43.13-1B, sección 4-59. Material: acero resistente a la corrosión tipo 301, de un cuarto de dureza, 0,016~in.
\item \textbf{Panel de quilla:} su reemplazo debe hacerse en fábrica, en una plantilla. Las reparaciones pueden hacerse según la AC 43.13-1B, secciones 4-58 y 4-59. Material: chapa de aleación de aluminio 2024-T3 \emph{clad}, 0,025~in.
\item \textbf{Estructura de los asientos:} para preservar la resistencia al impacto, las reparaciones se limitan a reemplazar los componentes dañados.
\end{itemize}

\subsection{Carenados, capós y paneles}

\begin{precaucion}
Los \textbf{faldones laterales} son un conjunto de paneles estructurales y deben estar instalados para la prueba de motor y el vuelo. El \textbf{carenado del mástil} y el \textbf{capó trasero} deben estar instalados para el vuelo. Todos los paneles de inspección de la cabina deben estar instalados para el vuelo; salvo los faldones laterales, pueden dejarse sin instalar para la prueba de motor.
\end{precaucion}

La costilla superior del carenado del mástil va montada en el conjunto del tubo del plato oscilante; la inferior está abrazada al mástil de la caja principal. El tubo pitot está en la parte delantera inferior del carenado. El respiradero de combustible pasa por la costilla inferior y debe tener unos 0,30~in de holgura con el capó trasero. El capó trasero es de dos piezas de chapa de aluminio, con un guardalluvia removible del cono de cola.

El respaldo del asiento izquierdo lleva la válvula de combustible; los cojines del respaldo se fijan con remaches ciegos. Los paneles de inspección de los controles del cíclico son cuatro: panel central del vientre, dos horizontales y uno vertical en el centro de la cabina, la consola de instrumentos y el panel del tubo transversal delantero.

\section{Bastidores de tubos de acero soldados}

\begin{advertencia}
Todas las estructuras de tubos de acero soldados del helicóptero tienen alivio de tensiones. No se permiten reparaciones por soldadura fuera de Robinson Helicopter Co.
\end{advertencia}

\begin{table}[htbp]
\centering\small
\caption{Bastidores según el modelo.}
\begin{tabularx}{\linewidth}{@{}lY@{}}
\toprule
Modelo & Bastidores \\
\midrule
Alpha y Beta & Superior A020-2 (S/N 0400 y posteriores); inferior izquierdo A046-1; inferior derecho A046-2; montante derecho A046-3 \\
Standard y HP & Superior A020-2; inferior izquierdo A020-1; inferior derecho A020-84; montante derecho A020-3 \\
Mariner (opcional) & Superior A047-1 con puntos de amarre \\
\bottomrule
\end{tabularx}
\end{table}

\subsection{Retiro e instalación de los bastidores}

Para retirar el \textbf{bastidor izquierdo} se retiran las palas, el embrague, la caja principal, el cono de cola, el grupo motor, los respaldos y paneles centrales; luego se desconectan los tres puntos inferiores delanteros del cortafuego y el superior trasero izquierdo, el perno trasero NAS1306 del soporte del tren, los dos puntos superiores en el cortafuego vertical y los pernos pasantes al bastidor superior. Valores clave de la instalación:

\begin{checklist}
\chk{Perno de llave interior bastidor superior--izquierdo (a través del cortafuego horizontal)}{75 ft-lb (húmedo)}
\chk{Pernos NAS1304 del soporte superior al cortafuego vertical}{100 in-lb + arrastre de tuerca}
\chk{Perno NAS1306 de unión del tren al soporte del bastidor}{300 in-lb + arrastre; \emph{palnuts}}
\chk{Perno NAS1304 del soporte superior trasero}{100 in-lb + arrastre; \emph{palnut}}
\end{checklist}

El soporte inferior en el fuselaje lleva una arandela de gran área AN970-4 bajo la cabeza del perno, junto a la oreja del bastidor; el soporte exterior requiere una arandela entre el bastidor y el revestimiento.

El \textbf{bastidor derecho} se retira de forma similar (cuatro puntos delanteros del cortafuego y dos pernos pasantes). En la instalación, el montante trasero entre el bastidor inferior derecho y el superior se aprieta según la tabla de torques y lleva \emph{palnuts}; selle el diámetro interior de los agujeros del montante con sellador B270-1, sin sellador en las roscas.

\begin{precaucion}
Antes de desconectar el \textbf{bastidor superior}, el grupo motor debe retirarse o sostenerse: si no se sostiene, o si el soporte se desplaza, se dañarán gravemente el cortafuego y los bastidores inferiores.
\end{precaucion}

Instalación del bastidor superior:
\begin{pasos}
\item Limpie el bastidor y los puntos de unión de sellador, grasa y aceite viejos.
\item Instale los pernos de llave interior con antiagarrotante, sin apretar, y los dos pernos pasantes delanteros NAS1306 (NAS6606).
\item Apriete: NAS1306 (NAS6606) a 300~in-lb; NAS1351-8H40P a 75~ft-lb (húmedo).
\item Mida la separación en el soporte delantero y suplemente si es de \textbf{0,032~in o más}. Si se reinstala el mismo bastidor, reponga los suplementos en su posición original.
\item Instale los cuatro pernos del soporte delantero (sin arandelas bajo las cabezas de los dos pernos B238-1 del túnel de los tubos) y apriételos a 40~in-lb + arrastre.
\item Instale los pernos NAS1304 (NAS6604) del soporte central, la unión trasera al cortafuego horizontal y la unión trasera al bastidor inferior, con su torque.
\item Selle los cortafuegos en los puntos de unión con Coast Pro-Seal 890B2, para que no se filtre combustible si hay una fuga de un tanque.
\end{pasos}

\begin{nota}
Los bastidores superiores A020-2 S/N 0002 a 0399 requieren la abrazadera A960-1 y la correa A961-1 según el R22 SB-26.
\end{nota}

\textbf{Montante vertical:} el extremo inferior va en la cara trasera de la oreja del bastidor inferior. Instale los pernos NAS6604-3 húmedos con sellador B270-1 en los vástagos (roscas limpias y secas), con torque especial (\S\,23-33), \emph{palnuts} con torque estándar y franja de torque.

\section{Cono de cola}

\begin{precaucion}
Si el cono de cola tiene el soporte B379-1 o -3, deben estar instalados (o incorporarse) los balancines A958-1 (equipo original desde el S/N 0357) y debe instalarse el estabilizador horizontal B375-2. Un cono con soporte B379 no puede instalarse en un Mariner. El estabilizador A044-1 no puede instalarse en un cono con soporte B379, y el soporte de estabilizadores verticales B902-1 o -2 solo puede instalarse en un cono que lo tenga.
\end{precaucion}

\textbf{Retiro:} abra los disyuntores de las luces y antenas del cono; retire el carenado A706-1; desconecte el cableado y las antenas; retire la ferretería entre la horquilla delantera del eje de cola y la placa flexible intermedia A947-2 (sostenga el eje con un bloque de espuma bien visible); desconecte el tubo A121-17 del balancín intermedio A331-1; sostenga el cono y retire la ferretería de unión al bastidor superior.

\textbf{Instalación:} el eje de cola, la caja de cola y el tubo A121-17 se instalan en el cono antes de montarlo. Revise el interior (sin residuos; los bujes de los mamparos impiden el contacto del tubo con el metal; el cableado tiene pasacables; el amortiguador está bien orientado). Posicione el cono sin pellizcar cables, instale la ferretería con torque estándar y franjas; conecte el tubo al balancín; instale y suplemente la placa flexible intermedia (\S\,7.330); mida el descentramiento del eje (\S\,7.340); conecte y pruebe la luz de posición, el estrobo y la luz \textsf{TR CHIP}; verifique las holguras con el bastidor; instale el estabilizador y el empenaje.

\subsection{Reemplazo del cono de cola (montaje y perforado)}

\begin{pasos}
\item Nivele el helicóptero longitudinal y lateralmente (\S\,18-10). Verifique que el número de parte del cono corresponda al modelo; deslice el cono sobre el bastidor y apoye su extremo trasero en un soporte.
\item \textbf{Nivelación lateral:} introduzca dos pernos en los agujeros de montaje de la caja de cola del mamparo A148 y ponga un nivel sobre sus vástagos; gire el cono hasta nivelarlos.
\item \textbf{Alineación lateral:} mida a cada lado, desde los pernos de unión exteriores del bastidor hasta la base del nivel en el mamparo A148; mueva el extremo trasero hasta igualar ambas medidas.
\item \textbf{Alineación vertical:} con un nivel de agua, la diferencia de altura entre el borde inferior delantero del cono y el eje de los pernos izquierdo y derecho en el mamparo A148 debe ser \textbf{16,21 $\pm$ 0,20~in}. (Por las tolerancias, la altura tiende a disminuir al retirar el soporte.)
\item \textbf{Distancias al borde:} holgura mínima de 0,010~in entre los tubos del bastidor y el borde delantero del cono (4 lugares), y mínimo 0,365~in entre una regla apoyada en la cara delantera del cono y el vástago del perno insertado en el resalte del bastidor (4 lugares).
\item \textbf{Perforado} (proteja las correas de las virutas): fije el cono con prensas en tres puntos (no en la oreja izquierda); retrotaladre desde adentro con broca de centro de vástago 0,250~in y luego broca helicoidal de 0,250~in; asegure con el sujetador y repita en cada punto. Taladre la oreja izquierda a través del cono protegiendo el tubo con una chapa.
\item Retire el cono y desbarbe. Centre los dobladores B288-1 o -2 sobre cada agujero, enrasados con el borde delantero (mínimo 0,345~in según la figura del MM; los remaches no deben perforar el radio del mamparo B070-1); taládrelos a 0,250~in, fíjelos con pernos NAS1304-3 y taladre los agujeros piloto con broca \#30; agregue los dobladores B288-3 encima y repita. Remache con MS20470AD4, lime al ras si sobresale, selle con B270-1, imprima (cromato de zinc o epóxico) y pinte.
\end{pasos}

\subsection{Límites de daño del cono de cola}

Las reparaciones se limitan a eliminar rayas dentro de límites y repintar. Si se exceden los límites, reemplace el cono o envíelo a RHC.

\begin{checklist}
\chk{Raya con eje a más de 15° de la línea central del cono}{0,005 in de profundidad máx.}
\chk{Raya con eje a menos de 15° de la línea central}{0,010 in de profundidad máx.}
\chk{Radio de transición mínimo al pulir rayas}{0,10 in}
\chk{Abolladura: fondo liso y redondeado, sin mellas ni grietas}{radio mín. 0,125 in}
\chk{Profundidad máxima de abolladura}{0,030 in}
\chk{Diámetro máximo de abolladura}{1,250 in}
\chk{Abolladuras por estación (pulgada) del cono}{una}
\chk{Distancia mínima entre estaciones con abolladuras}{4,000 in}
\end{checklist}

\begin{nota}
El radio mínimo de 0,125~in puede verificarse con una bola de rodamiento de 0,250~in de diámetro: póngala en la abolladura y retroilumine; si se ve luz entre la bola y el revestimiento (la bola no toca el fondo), el radio es menor que 0,125~in.
\end{nota}

\section{Estabilizadores y empenaje}

\subsection{Estabilizador horizontal B375-2}

Sostenga el estabilizador durante el retiro y la instalación cuando se retiran los pernos superiores, para no doblar el soporte B379 remachado al cono. Con el estabilizador retirado se recomienda instalar los soportes MT023-1 con la ferretería existente. Al instalarlo, aplique torque especial (\S\,23-33) y franja de torque. El B375-2 no puede instalarse en un Mariner.

\subsection{Empenaje}

Al instalar el empenaje en el mamparo A148 use dos pernos \textbf{NAS6604-28} si no se instalará el lastre de empenaje A301-5, o \textbf{NAS6604-44} si se instalará (\S\,18-32), con torque estándar y franjas. Reconecte la luz de posición, fije los cables con la abrazadera y amarras MS3367-4-9 (sin apretar en exceso, cortadas al ras) y pruebe la luz de posición y la \textsf{TR CHIP}.

\textbf{Reemplazo del empenaje A004-2:} nivele la aeronave; presente el empenaje con dos pernos NAS1304-28 en los agujeros superiores del A044-1 (puede usarse una prensa en C con tacos de madera). Ponga la herramienta de reglaje MT068-1 sobre el estabilizador horizontal, a 3,25--3,75~in del borde exterior del vertical, con el transportador o nivel sobre la superficie plana \emph{delantera} de la herramienta si el cono es A023-1 o -22 (Standard, HP y Mariner) o \emph{trasera} si es A023-20 o -23 (Alpha y Beta). Ajuste el ángulo de incidencia empujando o tirando levemente del patín de cola hasta nivelar esa superficie, y nivele lateralmente el borde de fuga a \textbf{0° $\pm$ 0,5°}. Taladre a 0,250~in usando los agujeros existentes como guía (primero el trasero, verificando de nuevo, y luego el delantero), desbarbe e instale.

\subsection{Estabilizadores verticales}

Al instalar los estabilizadores verticales superior A042-1 e inferior A043-1 sobre el soporte B902 (o sobre el A044-1), verifique una separación de \textbf{0,030 a 0,120~in} entre los bordes de su revestimiento y el revestimiento del soporte o estabilizador; lime los bordes si es necesario y aplique conversión química e imprimante al aluminio desnudo. Pernos con torque especial y franja. Al reemplazarlos, taladre las grapas A554 con broca de 0,144~in.

\subsection{Reparación del estabilizador horizontal A044-1}

Se permite una sola abolladura en el borde de ataque, por fuera de los estabilizadores verticales, si: (1) no tiene más de 0,050~in de profundidad; (2) tiene fondo liso con radio mínimo de 0,125~in, sin mellas ni grietas; (3) mide menos de 1,25~in en envergadura. Puede eliminarse por conformado. El reemplazo del revestimiento, los daños a los largueros y a las uniones delantera o media de los verticales no son reparables en campo. Solo la costilla de punta exterior B722-3 puede reemplazarse en campo (retirando los remaches NAS1919B04S01); todo lo demás requiere la plantilla de fábrica.

\subsection{Patín de cola}

Se fija al estabilizador vertical inferior A043-1 con ferretería de torque estándar y franja.

\section{Tren de aterrizaje}

\subsection{Descripción}

Tren de patines. La mayoría de los aterrizajes duros se absorben elásticamente; en uno extremadamente duro, los montantes se abren hacia arriba y afuera al ceder el tubo transversal. Una leve deformación es aceptable. \textbf{Reemplace el tubo transversal delantero si su deformación lo lleva a tocar su cubierta.} Una deformación que deje el patín de cola a menos de \textbf{34~in} del suelo (\textbf{24~in} en el Standard o HP; \textbf{37~in} en el Mariner, medido desde el estabilizador del flotador), con el helicóptero vacío sobre superficie dura y nivelada, exige reemplazar el o los tubos transversales. Los Mariner tienen flotadores utilitarios inflados permanentemente, montados en los patines, y extensiones traseras de patín que los sostienen.

\subsection{Retiro e instalación del tren}

\begin{pasos}
\item Retire la cubierta del tubo transversal delantero y levante el helicóptero con gatas o izaje (capítulo~\ref{cap:mm-datos-servicio}).
\item Desconecte el cable de masa A936-2 en el montante trasero derecho, si existe.
\item Retire los dos pernos de montaje izquierdos y apoye ese lado en el suelo (puede tirar del tren a la derecha para pivotar el grillete delantero derecho y acceder a su perno); luego los dos derechos.
\item Si el helicóptero se levantó con el bloque de contorno, retire un patín para sacar el tren.
\end{pasos}

\begin{nota}
Al reemplazar un tren rev.~AH o anterior por uno rev.~AI o posterior, reemplace también los dos soportes A014-7 por soportes A574-2. Todos los pernos de montaje se instalan desde atrás. En los trenes rev.~AH o anteriores, ponga arandelas de diámetro grande a ambos lados de cada montante para retener los espaciadores A278-1.
\end{nota}

Al instalar, monte primero el lado derecho y luego el izquierdo (deflectando con cuidado el bastidor inferior para alinear), con torque estándar y franjas. Si no se instala el cable de masa A936-2, se requieren los dos soportes A574-2 para la puesta a masa. \textbf{Verifique que la resistencia entre cada soporte de rueda de maniobra y la estructura de la cabina sea menor que 5~ohm.}

\subsection{Tubos transversales}

Reemplácelos si la altura del patín de cola cae por debajo de los valores indicados (vea también la inspección de aterrizaje duro en el capítulo~\ref{cap:mm-inspecciones}).

\begin{itemize}
\item \textbf{A241-3 con soldaduras de montante A239-4:} saque los montantes golpeando con martillo de cara blanda; los espaciadores se sacan con un botador de bronce o aluminio. Al instalar, aplique imprimante aprobado en el diámetro exterior de los espaciadores A243-6 y en las superficies de contacto con los montantes; con el imprimante húmedo, arme y alinee; instale los pernos a mano y selle todas las uniones con imprimante.
\item \textbf{A241-2 con montantes A014-18:} saque los espaciadores con un botador y retire los montantes. Si un montante está pegado: instale el espaciador del lado opuesto, suelte el montante pegado del patín, apoye el tubo cerca de él y gírelo en ambos sentidos para romper la adherencia de la pintura. Al instalar, imprimante húmedo en las superficies de contacto y en los espaciadores A278-1, y sellado de las uniones.
\end{itemize}

\subsection{Tubos de patín}

\textbf{Reemplace los patines} con rayas longitudinales de más de \textbf{0,030~in} de profundidad (sin límite de largo), rayas circunferenciales de más de \textbf{0,030~in} de profundidad o \textbf{0,50~in} de largo, o una zona plana de más de \textbf{0,50~in} de ancho. Para retirarlo, levante un lado con una gata bajo el extremo del tubo transversal trasero, 1~in hacia adentro del herraje, y retire los 4 pernos de cada montante. Al instalar, torque estándar (hay tuercas remachadas autoblocantes dentro del patín).

\subsection{Zapatas de los patines}

\textbf{Reemplace las zapatas A667-5 y -6 de menos de 0,050~in de espesor.}

\begin{pasos}
\item Instale las ruedas de maniobra. Para las zapatas delanteras e intermedias, empuje la cola hacia abajo y ponga un taco de madera de 4--6~in unas 4~in por delante de las zapatas intermedias.
\item Retire los pernos de las zapatas delanteras e intermedias y descarte zapatas y tornillos; instale los nuevos con torque estándar. (Algunos R22 antiguos tenían las zapatas delanteras fijadas con remaches ciegos de aluminio MS20600AD6-4: taladre la cabeza y saque el vástago.)
\item Para las zapatas traseras, incline el helicóptero hasta que el patín de cola toque el suelo, ponga un taco de 4--6~in lo más atrás posible bajo cada patín, presione las puntas de los patines o levante el cono hasta que el helicóptero se incline hacia adelante, y lastre cada punta con unas 25~lb (o apoye el patín de cola). Reemplace las zapatas igual que las delanteras.
\end{pasos}

\textbf{Soporte de rueda de maniobra:} dos pernos de unión al patín, apretados a 50~in-lb + arrastre de tuerca.

\subsection{Mantenimiento preventivo del tren}

\begin{itemize}
\item Mantenga despejados los agujeros de drenaje bajo los patines (dos por patín, justo detrás de las dos zapatas delanteras).
\item Retoque la pintura de patines, tubos transversales y montantes para prevenir la corrosión, con materiales aprobados.
\item Con uso intensivo, inspeccione las zapatas en cada prevuelo.
\item Reinstale con adhesivo estructural las tapas de lluvia perdidas o sueltas.
\end{itemize}

\section{Tren de flotadores (Mariner)}

El R22 Mariner puede volar con o sin tren de flotadores. Para retirar los flotadores debe instalarse un tren estándar completo; no se separan los flotadores de los patines.

\begin{precaucion}
El tren de flotadores solo puede instalarse en helicópteros R22 Mariner, que tienen el sufijo ``M'' en el número de serie de la placa de datos.
\end{precaucion}

\textbf{Cambio a tren estándar:} retire el tren de flotadores completo e instale el estándar; retire el estabilizador de flotador del estabilizador vertical inferior e instale el patín de cola estándar (pernos a 30~in-lb); recalcule el peso y balance (tabla~\ref{tab:pyb-flotadores}); revise las limitaciones sin flotadores en el suplemento Mariner; haga prevuelo, prueba de motor y vuelo de chequeo, y \textbf{ajuste la RPM de autorrotación} (\S\,10.250; normalmente hay que alargar las varillas de paso de ½ a 1 vuelta).

\textbf{Cambio a tren de flotadores:} el procedimiento inverso, con la ferretería existente salvo las \emph{palnuts}, que deben reemplazarse; el estabilizador de flotador reemplaza al patín (pernos a 30~in-lb); recalcule el peso y balance y ajuste la RPM de autorrotación (normalmente acortando las varillas de ½ a 1 vuelta). Para la presión de los flotadores, vea el suplemento Mariner del POH.

\begin{table}[htbp]
\centering\small
\caption{Datos de peso y balance para el cambio de tren (Mariner).}
\label{tab:pyb-flotadores}
\begin{tabularx}{\linewidth}{@{}lYrrr@{}}
\toprule
P/N & Descripción & Peso (lb) & Brazo (in) & Momento (in-lb) \\
\midrule
A034-1 & Tren con flotadores & +62,9 & 94,6 & +5950 \\
A014-1 & Tren sin flotadores & $-32,5$ & 85,5 & $-2779$ \\
A970 y A971 & Estabilizador de flotador en la aleta inferior & +0,6 & 270 & +162 \\
A470-1 & Patín de cola estándar & $-0,2$ & 272,5 & $-54$ \\
\bottomrule
\end{tabularx}
\end{table}

\subsection{Flotadores}

\textbf{Retiro:} desate y retire los cordones en cada montante y desabroche las correas de retención; retire los pernos y tuercas que unen el flotador al patín y a la extensión (están sellados y pueden costar); retire las tiras retenedoras y separe con cuidado la faja (\emph{girt}) de la extensión (hay sellador entre faja y patín en cada agujero).

\textbf{Instalación:} retire el sellador viejo. Presente el flotador, alinee los agujeros, coloque las tiras retenedoras con pernos temporales y, una sección de faja por vez, selle alrededor de los agujeros entre faja y patín con sellador 890-B2; instale los pernos del largo correcto (2 a 4 hilos expuestos), sellando bajo cabezas y tuercas, y apriete los NAS6604 o NAS1304 a \textbf{100~in-lb} + arrastre. Imprima y pinte de negro los pernos y tuercas expuestos. Ate el flotador al montante con un patrón cruzado, pase las correas por los bolsillos exteriores, rodee el flotador por debajo (pasando sobre el patín), dé una vuelta al montante, abroche, tense y guarde el sobrante en el bolsillo para que no flamee en vuelo. Infle lentamente a la presión de operación (de a 0,5~psig por cámara si se hace sin múltiple). Para reparaciones, siga las instrucciones del kit de reparación de flotadores.
